\BourbakiBackMatter{参考文献}

{[}1{]} E. ARTIN -- ``Zur Theorie der hypercomplexen Zahlen'', Abh. math. Sem. Univ. Hamburg 5 (1928), p.~251--260; Collected papers, Reading (Addison Wesley), 1965, p.~307--316.

{[}2{]} R. BRAUER -- ``Über Systeme hypercomplexer Zahlen'', Math. Zeitschr. 30 (1929), p.~79--107; Collected papers, vol.~I, Cambridge, Massachusetts (The MIT Press), 1980, p.~40--68.

{[}3{]} W. BURNSIDE - ``论任意线性替换群的可约性条件'', Proc. Lond. Math. Soc. 3 (1905), p.~430-434.

{[}4{]} É. CARTAN -- ``Les groupes bilinéaires et les systèmes de nombres complexes'', Ann. Fac. Sc. Toulouse 12 (1898), p.~1--99; Œuvres complètes, partie II, vol.~I, Paris (Gauthier-Villars), 1952, p.~7--105.

{[}5{]} \_\_\_\_ ``Nombres complexes'', Encycl. Sci. Math., éd. française 15 (1908); Œuvres complètes, partie II, vol.~I, Paris (Gauthier-Villars), 1952, p.~107--246.

{[}6{]} A. CAYLEY -- ``论代数偶'', Phil. Mag. (3) 27 (1845), p.~38--40; Collected mathematical papers, t. I, Cambridge, 1889, p.~128--131.

{[}7{]} \_\_\_\_ ``论依赖于方程 \(\theta^{n}=1\) 的群论'', Phil. Mag. (4) 7 (1854), p.~40--47; Collected mathematical papers, t. II, Cambridge, 1889, p.~123--130.

{[}8{]} \_\_\_\_``矩阵理论专论'', Phil. Trans. R. Soc. Lond. 148 (1858), p.~17--37; Collected mathematical papers, t. II, Cambridge, 1889, p.~475--496.

{[}9{]} W. CLIFFORD -- ``双四元数初步概要'', Proc. Lond. Math. Soc. 4 (1873), p.~381--395; Mathematical papers, London (Macmillan), 1882, p.~181--200.

{[}10{]} \_\_\_\_ ``论几何代数的分类'', (1877), Mathematical papers, London (Macmillan), 1882, p.~397--401.

{[}11{]} \_\_\_\_``格拉斯曼扩张代数的应用'', Amer. Journ. of Math. 1 (1878), p.~350--358; Mathematical papers, London (Macmillan), 1882, p.~266--276.

{[}12{]} R. DEDEKIND -- ``Zur Theorie der aus n Haupteinheiten gebildeten komplexen Grössen'', Gött. Nachr. (1885), p.~141--159; Gesammelte mathematische Werke, t. II, Braunschweig (Vieweg), 1931--32, p.~1--19.

{[}13{]} \_\_\_\_ ``Aus Briefen an Frobenius'', (1886), Gesammelte mathematische Werke, t. II, Braunschweig (Vieweg), 1931--32, p.~414--442.

{[}14{]} \_\_\_\_ ``Über eine Erweiterung des Symbols (a, b) in der Theorie der Moduln'', Gött. Nachr. (1895), p.~183--188; Gesammelte mathematische Werke, t. II, Braunschweig (Vieweg), 1931--32, p.~59--86.

{[}15{]} \_\_\_\_``Über Gruppen, deren sämtliche Teiler Normalteiler sind'', Math. Ann. 48 (1897), p.~548--561; Gesammelte mathematische Werke, t. II, Braunschweig (Vieweg), 1931--32, p.~87--102.

{[}16{]} M. DEURING - Algebren, Erg. der Math., vol.~4, Berlin (Springer), 1935.

{[}17{]} L. DICKSON - ``线性结合代数与阿贝尔方程'', Trans. Amer. Math. Soc. 15 (1914), p.~31-46; Mathematical papers, vol.~II, New York (Chelsea), 1975, p.~367-382.

{[}18{]} G. FROBENIUS -- ``Über lineare Substitutionen und bilineare Formen'', Crelle's J. 84 (1878), p.~1--63; Gesammelte Abhandlungen, Band I, Berlin (Springer), 1968, p.~343--405.

{[}19{]} \_\_\_\_ ``Über die Primfactoren der Gruppendeterminante'', Berliner Sitzungsber. (1896), p.~1343--1382; Gesammelte Abhandlungen, Band III, Berlin (Springer), 1968, p.~38--77.

{[}20{]} \_\_\_\_ ``Über Gruppencharaktere'', Berliner Sitzungsber. (1896), p.~985--1021; Gesammelte Abhandlungen, Band III, Berlin (Springer), 1968, p.~1--37.

{[}21{]} \_\_\_\_``Über die Darstellung der endlichen Gruppen durch lineare Substitutionen'', Berliner Sitzungsber. (1897), p.~944--1015; Gesammelte Abhandlungen, Band III, Berlin (Springer), 1968, p.~82--103.

{[}22{]} \_\_\_\_ ``Theorie der hypercomplexen Grössen'', Berliner Sitzungsber. (1903), p.~504--537 与 634--645; Gesammelte Abhandlungen, Band III, Berlin (Springer), 1968, p.~284--329.

{[}23{]} H. GRASSMANN -- ``Sur les différents genres de multiplication'', Crelle's J. 49 (1855), p.~123--141; Gesammelte math. und phys. Werke, t. II, Leipzig (Teubner), 1904, p.~199--217.

{[}24{]} \_\_\_\_ ``Die Ausdehnungslehre von 1862'', (1862), Gesammelte math. und phys. Werke, t. I, Leipzig (Teubner), 1896, p.~1--383.

{[}25{]} W. HAMILTON - 四元数讲义, Dublin (Hodges and Smith), 1853.

{[}26{]} O. HÖLDER -- ``Zurückführung einer beliebigen algebraischen Gleichung auf eine Kette von Gleichungen'', Math. Ann. 34 (1889), p.~26--56.

{[}27{]} C. HOPKINS - ``具有左理想极小条件的环'', Ann. of Math. 40 (1939), p.~712-730.

{[}28{]} W. KRULL -- ``Über verallgemeinerte endliche Abelsche Gruppen'', Math. Zeitschr. 23 (1925), p.~161--196; Gesammelte Abhandlungen, vol.~1, Berlin (de Gruyter), 1999, p.~263--298.

{[}29{]} \_\_\_\_ ``Theorie und Anwendung der verallgemeinerten Abelschen Gruppen'', Sitzungsber. Heidelberger Akad. (1926), p.~1--32; Gesammelte Abhandlungen, vol.~1, Berlin (de Gruyter), 1999, p.~299--328.

{[}30{]} E. LAGUERRE -- ``Sur le calcul des systèmes linéaires'', J. de l'école polytechnique 42 (1867), p.~215--264; Œuvres, t. I, Paris (Gauthier-Villars).

{[}31{]} T. MOLIEN - ``Ueber Systeme höherer complexer Zahlen'', Math. Ann. 41 (1893), p.~83-156.

{[}32{]} \_\_\_\_``Ueber die Invarianten der linearen Substitutionsgruppen'', Berliner Sitzungsber. (1897), p.~1152--1156.

{[}33{]} E. NOETHER -- ``Abstrakter Aufbau der Idealtheorie in algebraischen Zahl- und Funktionenkörpern'', Math. Ann. 96 (1927), p.~26--61; Gesammelte Abhandlungen, Berlin (Springer), 1983, p.~493--528.

{[}34{]} \_\_\_\_ ``Hyperkomplexe Grössen und Darstellungstheorie'', Math. Zeitschr. 30 (1929), p.~641--692; Gesammelte Abhandlungen, Berlin (Springer), 1983, p.~563--614.

{[}35{]} \_\_\_\_ ``Nichtkommutative Algebren'', Math. Zeitschr. 37 (1933), p.~514--541; Gesammelte Abhandlungen, Berlin (Springer), 1983, p.~642--669.

{[}36{]} E. NOETHER et R. BRAUER -- ``Über minimale Zerfällungskörper irreduzibler Darstellungen'', Berliner Sitzungsber. (1927), p.~221--228; Emmy Noether Gesammelte Abhandlungen, Berlin (Springer), 1983, p.~552--559.

{[}37{]} E. NOETHER et W. SCHMEIDLER -- ``Moduln in nichtkommutativen Bereichen'', Math. Zeitschr. 8 (1920), p.~1--35; Emmy Noether Gesammelte Abhandlungen, Berlin (Springer), 1983, p.~318--352.

{[}38{]} B. PEIRCE - ``线性结合代数'', Amer. Journ. of Math. 4 (1881), p.~97-221.

{[}39{]} C. PEIRCE - ``论除法无歧义的代数'', Amer. Journ. of Math. 4 (1881), p.~225-229.

{[}40{]} \_\_\_\_ ``论代数的相对形式'', Amer. Journ. of Math. 4 (1881), p.~221--225.

{[}41{]} H. POINCARÉ -- ``Sur les nombres complexes'', C. R. Acad. Sci. 99 (1884), p.~740--742; Œuvres, t. V, Paris (Gauthier-Villars), 1916--1954, p.~77--79.

{[}42{]} \_\_\_\_ ``Sur l'intégration algébrique des équations linéaires et les périodes des intégrales abéliennes'', Journ. de Math. (5) 9 (1903), p.~139--212; Œuvres, t. III, Paris (Gauthier-Villars), 1916--1954, p.~106--166.

{[}43{]} G. SCHEFFERS -- ``Zurückführung complexer Zahlensysteme auf typische Formen'', Math. Ann. 39 (1891), p.~293--390.

{[}44{]} I. SCHUR -- ``Über eine Klasse von Matrizen, die sich einer gegebenen Matrix zuordnen lassen'', Diss., Berlin, 1901; Gesammelte Abhandlungen, Band I, Berlin (Springer), 1973, p.~1--72.

{[}45{]} \_\_\_\_``Über die Darstellung der endlichen Gruppen durch gebrochene lineare Substitutionen'', Crelle's J. 127 (1904), p.~20--50; Gesammelte Abhandlungen, Band I, Berlin (Springer), 1973, p.~86--116.

{[}46{]} \_\_\_\_ ``Neue Begründung der Theorie der Gruppencharaktere'', Berliner Sitzungsber. (1905), p.~406--432; Gesammelte Abhandlungen, Band I, Berlin (Springer), 1973, p.~143--169.

{[}47{]} \_\_\_\_``Arithmetische Untersuchungen über endliche Gruppen linearer Substitutionen'', Berliner Sitzungsber. (1906), p.~164--184; Gesammelte Abhandlungen, Band I, Berlin (Springer), 1973, p.~177--197.

{[}48{]} J.-A. DE SÉGUIER - Éléments de la théorie des groupes abstraits, Paris (Gauthier-Villars), 1904.

{[}49{]} T. SKOLEM - Zur Theorie des assoziativen Zahlensysteme, Oslo Vid. Akad. Skr. I, 1927.

{[}50{]} J. M. WEDDERBURN - ``关于有限代数的一个定理'', Trans. Amer. Math. Soc. 6 (1905), p.~349-352.

{[}51{]} \_\_\_\_``论超复数'', Proc. Lond. Math. Soc. (2) 6 (1908), p.~77--118.

{[}52{]} \_\_\_\_``一类本原代数'', Trans. Amer. Math. Soc. 15 (1914), p.~162--166.

{[}53{]} \_\_\_\_ ``论可除代数'', Trans. Amer. Math. Soc. 22 (1921), p.~129--135.

{[}54{]} K. WEIERSTRASS -- ``Zur Theorie der aus n Haupteinheiten gebildeten complexen Grössen'', Gött. Nachr. (1884), p.~395--414; Mathematische Werke, t. II, Berlin (Mayer und Müller), 1895, p.~311--332.

